\manualchapter{Response and troubleshooting}{sec:operation}
\subsection{Host-rate timing and quantized level}
The envelope advances once per host sample rather than running at a fixed
32\,kHz chip rate. Changing Rack's sample rate changes its timing. Release
subtracts eight internal units per sample from a maximum of 2048, giving at
most 256 samples (about 5.33\,ms at 48\,kHz). Output quantization can reach
zero sooner. There is no adjustable release-rate control in the current path.

The output passes through signed 8-bit scaling and is approximately
$\pm10$\,V at full amplitude. Small amplitude values expose coarse steps and
can shorten the audible tail. For smooth fades, use another envelope or
smooth the result externally rather than expecting continuous chip levels.

\subsection{Two related contours}
Send the same gate to both lanes. Set lane 1 to a fast attack and lane 2 to
a slower attack, then use them for amplitude and filter modulation. The
second lane needs its own gate cable; lanes are not normalled together.
Use a clock on RETRIG while holding GATE to repeat the attack.

Both lanes use the widest gate/retrigger input's channel count, from 1--16.
Each socket reads matching cable channels without mono-to-poly repetition.
Panel settings are shared across those channels and saved by Rack; envelope
progress is not serialized. Reset restores parameters, but this module has
no extra envelope-state reset handler or custom context option.

If there is no contour, check Amplitude and GATE, then remove RETRIG while
diagnosing. If the tail is too short, the RR-marked slider cannot change the
fixed release. A reported Contour behavior bug remains tracked separately
in issue \href{https://github.com/Kautenja/PotatoChips/issues/98}{\#98}; this
manual documents the current path and does not claim that bug is fixed.
